\section{Worked Particle Models}
\label{sec:worked_particle_models}

The purpose of this section is to apply the variational method as a complete calculation. For each model, we identify the coordinate, write the Lagrangian, derive the Euler--Lagrange equation, solve or simplify the equation, and interpret the result.

The examples are deliberately standard. Their value lies not in novelty but in making the procedure operational.

\subsection{Model 1: the free particle}

Consider a particle of mass \(m\) moving on a line with no potential energy. The generalized coordinate is \(q(t)\), and the Lagrangian is

\[
L(q,\dot q)
=
\frac{m}{2}\dot q^{2}.
\]

The required derivatives are

\[
\frac{\partial L}{\partial q}=0,
\]

and

\[
\frac{\partial L}{\partial\dot q}
=
m\dot q.
\]

Therefore

\[
\frac{\mathrm{d}}{\mathrm{d}t}
\left(
\frac{\partial L}{\partial\dot q}
\right)
=
m\ddot q.
\]

The Euler--Lagrange equation gives

\[
m\ddot q=0,
\]

or

\[
\boxed{
\ddot q=0
}.
\]

Integrating once,

\[
\dot q=v_{0},
\]

where \(v_{0}\) is constant. Integrating again,

\[
\boxed{
q(t)=q_{0}+v_{0}t
}.
\]

The canonical momentum is

\[
p
=
\frac{\partial L}{\partial\dot q}
=
m\dot q,
\]

and the equation of motion states

\[
\dot p=0.
\]

Thus free motion corresponds to constant momentum and constant velocity.

\subsubsection{Fixed-endpoint solution}

Suppose

\[
q(t_{1})=q_{1},
\qquad
q(t_{2})=q_{2}.
\]

The solution satisfying these conditions is

\[
\boxed{
q(t)
=
q_{1}
+
\frac{q_{2}-q_{1}}{t_{2}-t_{1}}
(t-t_{1})
}.
\]

The physical path is the straight line in the \((t,q)\)-plane connecting the two endpoint events.

\subsubsection{Check by substitution}

Differentiating twice,

\[
\dot q
=
\frac{q_{2}-q_{1}}{t_{2}-t_{1}},
\qquad
\ddot q=0.
\]

The trajectory therefore satisfies the Euler--Lagrange equation exactly.

\subsection{Model 2: a particle in a general potential}

Let

\[
L(q,\dot q)
=
\frac{m}{2}\dot q^{2}-V(q).
\]

Then

\[
\frac{\partial L}{\partial q}
=
-V'(q),
\]

and

\[
\frac{\partial L}{\partial\dot q}
=
m\dot q.
\]

The Euler--Lagrange equation is

\[
m\ddot q+V'(q)=0.
\]

Defining the force by

\[
F(q)
=
-V'(q),
\]

we obtain

\[
\boxed{
m\ddot q=F(q)
}.
\]

This is Newton's second law. The Lagrangian and Newtonian formulations therefore describe the same dynamics for conservative forces.

\subsubsection{Energy check}

Multiply the equation of motion by \(\dot q\):

\[
m\ddot q\dot q
+
V'(q)\dot q
=0.
\]

Using

\[
m\ddot q\dot q
=
\frac{\mathrm{d}}{\mathrm{d}t}
\left(
\frac{m}{2}\dot q^{2}
\right),
\]

and

\[
V'(q)\dot q
=
\frac{\mathrm{d}V}{\mathrm{d}t},
\]

we find

\[
\frac{\mathrm{d}}{\mathrm{d}t}
\left(
\frac{m}{2}\dot q^{2}+V(q)
\right)
=0.
\]

Thus

\[
\boxed{
E
=
\frac{m}{2}\dot q^{2}+V(q)
=
\text{constant}
}.
\]

This conservation law provides a useful independent check of solutions.

\subsection{Model 3: motion in a uniform gravitational field}

Choose the vertical coordinate \(y\), positive upward. The potential energy is

\[
V(y)=mgy.
\]

The Lagrangian is

\[
L(y,\dot y)
=
\frac{m}{2}\dot y^{2}-mgy.
\]

We calculate

\[
\frac{\partial L}{\partial y}
=
-mg,
\]

and

\[
\frac{\partial L}{\partial\dot y}
=
m\dot y.
\]

Therefore

\[
\frac{\mathrm{d}}{\mathrm{d}t}
\left(
\frac{\partial L}{\partial\dot y}
\right)
=
m\ddot y.
\]

The Euler--Lagrange equation gives

\[
m\ddot y+mg=0.
\]

Dividing by \(m\),

\[
\boxed{
\ddot y=-g
}.
\]

Integrating,

\[
\dot y(t)=v_{0}-gt,
\]

and

\[
\boxed{
y(t)=y_{0}+v_{0}t-rac{g}{2}t^{2}
}.
\]

\subsubsection{Physical interpretation}

The mass cancels from the equation of motion. In this idealized model, all particles have the same gravitational acceleration regardless of mass.

The conserved energy is

\[
E
=
\frac{m}{2}\dot y^{2}+mgy.
\]

Substituting the solution verifies that the decrease in kinetic energy during upward motion equals the increase in gravitational potential energy.

\subsection{Model 4: the harmonic oscillator}

Consider a particle attached to a spring with spring constant \(k\). The potential energy is

\[
V(q)
=
\frac{k}{2}q^{2}.
\]

The Lagrangian is

\[
L(q,\dot q)
=
\frac{m}{2}\dot q^{2}
-
\frac{k}{2}q^{2}.
\]

We have

\[
\frac{\partial L}{\partial q}
=
-kq,
\]

and

\[
\frac{\partial L}{\partial\dot q}
=
m\dot q.
\]

The Euler--Lagrange equation gives

\[
m\ddot q+kq=0.
\]

Define

\[
\omega^{2}
=
\frac{k}{m}.
\]

Then

\[
\boxed{
\ddot q+\omega^{2}q=0
}.
\]

The general solution is

\[
\boxed{
q(t)
=
A\cos(\omega t)
+
B\sin(\omega t)
}.
\]

Equivalently,

\[
q(t)
=
C\cos(\omega t-\varphi).
\]

\subsubsection{Verification}

Differentiating twice,

\[
\ddot q(t)
=
-\omega^{2}
\left[
A\cos(\omega t)
+
B\sin(\omega t)
\right]
=
-\omega^{2}q(t).
\]

Thus the proposed solution satisfies the equation.

\subsubsection{Energy}

The conserved energy is

\[
E
=
\frac{m}{2}\dot q^{2}
+
\frac{k}{2}q^{2}.
\]

For

\[
q(t)=C\cos(\omega t-\varphi),
\]

we have

\[
\dot q(t)
=
-C\omega\sin(\omega t-\varphi).
\]

Therefore

\[
\begin{aligned}
E
&=
\frac{m}{2}
C^{2}\omega^{2}
\sin^{2}(\omega t-\varphi)
+
\frac{k}{2}
C^{2}
\cos^{2}(\omega t-\varphi)\\
&=
\frac{kC^{2}}{2}
\left[
\sin^{2}(\omega t-\varphi)
+
\cos^{2}(\omega t-\varphi)
\right]\\
&=
\boxed{
\frac{kC^{2}}{2}
}.
\end{aligned}
\]

The total energy remains constant while kinetic and potential energies exchange periodically.

\subsection{Model 5: a particle in two-dimensional polar coordinates}

Consider a particle moving in a plane. Introduce polar coordinates

\[
x=r\cos\theta,
\qquad
y=r\sin\theta.
\]

The velocity components are

\[
\dot x
=
\dot r\cos\theta-r\dot\theta\sin\theta,
\]

and

\[
\dot y
=
\dot r\sin\theta+r\dot\theta\cos\theta.
\]

Squaring and adding,

\[
\dot x^{2}+\dot y^{2}
=
\dot r^{2}+r^{2}\dot\theta^{2}.
\]

For a central potential \(V(r)\), the Lagrangian is

\[
\boxed{
L(r,\theta,\dot r,\dot\theta)
=
\frac{m}{2}
\left(
\dot r^{2}+r^{2}\dot\theta^{2}
\right)
-
V(r)
}.
\]

\subsubsection{Angular equation}

The coordinate \(\theta\) does not appear explicitly in \(L\):

\[
\frac{\partial L}{\partial\theta}=0.
\]

The conjugate momentum is

\[
p_{\theta}
=
\frac{\partial L}{\partial\dot\theta}
=
mr^{2}\dot\theta.
\]

The Euler--Lagrange equation gives

\[
\frac{\mathrm{d}}{\mathrm{d}t}
\left(mr^{2}\dot\theta\right)
=0.
\]

Therefore

\[
\boxed{
mr^{2}\dot\theta
=\ell
=\text{constant}
}.
\]

The conserved quantity \(\ell\) is the angular momentum about the origin.

\subsubsection{Radial equation}

For the radial coordinate,

\[
\frac{\partial L}{\partial\dot r}
=
m\dot r,
\]

so

\[
\frac{\mathrm{d}}{\mathrm{d}t}
\left(
\frac{\partial L}{\partial\dot r}
\right)
=
m\ddot r.
\]

Also,

\[
\frac{\partial L}{\partial r}
=
mr\dot\theta^{2}-V'(r).
\]

The radial Euler--Lagrange equation is

\[
m\ddot r
-
mr\dot\theta^{2}
+
V'(r)
=0.
\]

Equivalently,

\[
\boxed{
m\ddot r
=
mr\dot\theta^{2}
-
V'(r)
}.
\]

The term

\[
mr\dot\theta^{2}
\]

arises from the coordinate dependence of the kinetic energy. It is not inserted as an additional force by hand.

\subsubsection{Effective potential}

Using

\[
\dot\theta
=
\frac{\ell}{mr^{2}},
\]

the energy becomes

\[
E
=
\frac{m}{2}\dot r^{2}
+
\frac{\ell^{2}}{2mr^{2}}
+
V(r).
\]

Define the effective potential

\[
\boxed{
V_{\mathrm{eff}}(r)
=
V(r)
+
\frac{\ell^{2}}{2mr^{2}}
}.
\]

Then the radial motion takes the one-dimensional energy form

\[
E
=
\frac{m}{2}\dot r^{2}
+
V_{\mathrm{eff}}(r).
\]

This reduction shows how a multidimensional problem can be simplified using a cyclic coordinate and its conserved momentum.

\subsection{Model 6: the simple pendulum and the small-angle limit}

For a pendulum of length \(\ell\) and mass \(m\), the Lagrangian is

\[
L
=
\frac{m\ell^{2}}{2}\dot\theta^{2}
-
mg\ell(1-\cos\theta).
\]

The equation of motion is

\[
\boxed{
\ddot\theta
+
\frac{g}{\ell}
\sin\theta
=0
}.
\]

This equation is nonlinear because of the sine function.

For small angular displacements,

\[
|\theta|\ll 1,
\]

we use

\[
\sin\theta
\approx
\theta.
\]

The equation becomes

\[
\ddot\theta
+
\frac{g}{\ell}\theta
=0.
\]

This is the harmonic-oscillator equation with

\[
\omega_{0}
=
\sqrt{\frac{g}{\ell}}.
\]

The approximate solution is

\[
\boxed{
\theta(t)
=
A\cos(\omega_{0}t)
+
B\sin(\omega_{0}t)
}.
\]

The corresponding small-amplitude period is

\[
\boxed{
T
=
2\pi
\sqrt{\frac{\ell}{g}}
}.
\]

This example illustrates linearization: a nonlinear equation is approximated near an equilibrium by retaining only the leading term in a series expansion.

\subsection{Comparing the models}

The examples reveal a common pattern.

\begin{enumerate}
    \item The free particle has a cyclic position coordinate and conserved linear momentum.
    \item A potential introduces a generalized force through \(-V'(q)\).
    \item A quadratic potential produces linear harmonic motion.
    \item Curvilinear coordinates make the kinetic energy coordinate dependent.
    \item A cyclic angular coordinate produces conserved angular momentum.
    \item A nonlinear system may reduce to a linear oscillator near equilibrium.
\end{enumerate}

These are not isolated tricks. The same structural ideas recur in field theory: kinetic terms, potentials, symmetries, conserved quantities, coordinate choices, and linearization around simple backgrounds.

\subsection{A complete checking procedure}

For every model derived from an action, one should check:

\begin{enumerate}
    \item Are all terms in the Lagrangian dimensionally consistent?
    \item Were the generalized coordinates and constraints identified correctly?
    \item Was the total time derivative in the Euler--Lagrange equation calculated fully?
    \item Does the sign of the force agree with the physical potential?
    \item Does a proposed solution satisfy the differential equation?
    \item Are conserved quantities constant along the solution?
    \item Does the result reduce correctly in a simple limit?
\end{enumerate}

These checks often reveal errors more quickly than repeating the entire derivation.

\subsection{Chapter summary}

This chapter developed the action principle first as a mathematical construction and then as a practical method of deriving dynamics.

The action is a functional of the complete trajectory:

\[
S[q]
=
\int_{t_{1}}^{t_{2}}
L(q,\dot q,t)
\,\mathrm{d}t.
\]

A varied path is written as

\[
q_{\varepsilon}(t)
=
q(t)+\varepsilon\eta(t),
\]

and fixed endpoint positions imply

\[
\eta(t_{1})=\eta(t_{2})=0.
\]

Integration by parts separates the variation into a boundary term and an interior term. Stationarity for arbitrary admissible deformations produces the Euler--Lagrange equation

\[
\boxed{
\frac{\mathrm{d}}{\mathrm{d}t}
\left(
\frac{\partial L}{\partial\dot q^{i}}
\right)
-
\frac{\partial L}{\partial q^{i}}
=0
}.
\]

Generalized coordinates incorporate constraints and adapt the description to the geometry of the system. The worked models showed how the same procedure produces free motion, uniform acceleration, harmonic oscillation, angular-momentum conservation, effective radial dynamics, and the small-angle pendulum.

The next chapter extends every part of this construction from particle trajectories to fields. The generalized coordinate \(q(t)\) becomes a field \(\phi(x)\), the Lagrangian becomes a Lagrangian density, the time integral becomes a spacetime integral, and ordinary integration by parts becomes spacetime integration by parts.